\subsection{Properties of Equality} 
\begin{displaybox}[Symmetric Property of Equality]
For any real numbers $a$ and $b$:
\[
a=b
\quad\mbox{if and only if}\quad
b=a
\]
\end{displaybox}
\begin{displaybox}[Transitive Property of Equality]
For any real numbers $a$, $b$, and $c$:
\[
\mbox{If }a=b\mbox{ and }b=c\mbox{ then }a=c
\]
\end{displaybox}
\begin{displaybox}[Additive Property of Equality]
For any real numbers $a$, $b$, and $c$:
\[
a=b
\quad\mbox{if and only if}\quad
a+c = b+c
\]
%This property means that we can get an ``equivalent equation'' by adding the same number to both sides. 
We will sometimes think of this as ``subtracting,'' but remember that subtracting is just adding the \makeblank[1.5in].
\end{displaybox}
\begin{displaybox}[Multiplicative Property of Equality]
For any real numbers $a$, $b$, and $c$, with $c\neq 0$:
\[
a=b
\quad\mbox{if and only if}\quad
ac=bc
\]
%This property means that we can get an ``equivalent equation'' by multiplying both sides by the same number, as long as it isn't zero. 
We will sometimes think of this as ``dividing,'' but remember that dividing is just multiplying by the \makeblank[1.5in].
\end{displaybox}
\begin{example}
Assume $2x+2=4x-6$. Manipulate the equation as follows:
\begin{lpic}{}{4}
\regtextleft{0}{-1}{Add $3$ to both sides};
\regtextleft{0}{-2}{Multiply both sides by $2$};
\regtextleft{0}{-3}{Subtract $3x$ from both sides};
\regtextleft{0}{-4}{Switch the order of the sides};
\end{lpic}
To check that these equations are all equivalent, we'll see that they're all true when $x=4$.
\begin{lpic}{}{3}
\end{lpic}
\end{example}